\documentclass[12pt]{article}
\usepackage[top=1in, bottom=1in, left=1in, right=1in]{geometry}
\usepackage{hyperref}
\usepackage{graphicx}
\usepackage{listings}
\usepackage{xcolor}
\usepackage{tocloft}
\usepackage{titlesec}
\usepackage{setspace}
\usepackage{caption}

% Line spacing for readability and professional report formatting
\onehalfspacing

% Code block styling
\definecolor{codegreen}{rgb}{0,0.6,0}
\definecolor{codegray}{rgb}{0.5,0.5,0.5}
\definecolor{codepurple}{rgb}{0.58,0,0.82}
\definecolor{backcolour}{rgb}{0.95,0.95,0.92}

\lstdefinestyle{mystyle}{
    backgroundcolor=\color{backcolour},   
    commentstyle=\color{codegreen},
    keywordstyle=\color{magenta},
    numberstyle=\tiny\color{codegray},
    stringstyle=\color{codepurple},
    basicstyle=\ttfamily\footnotesize,
    breakatwhitespace=false,         
    breaklines=true,                 
    captionpos=b,                    
    keepspaces=true,                 
    numbers=left,                    
    numbersep=5pt,                  
    showspaces=false,                
    showstringspaces=false,
    showtabs=false,                  
    tabsize=2
}
\lstset{style=mystyle}

\title{\textbf{EconBet Technical Report: Frontend Architecture, UX Design, \& Implementation}}
\author{EconBet Development Team}
\date{\today}

\begin{document}

\maketitle
\thispagestyle{empty}
\newpage

\tableofcontents
\newpage

\section{Executive Summary}
This comprehensive technical report details the ideation, architectural decisions, user experience (UX) philosophy, and technical implementation of the frontend for the EconBet Terminal. EconBet is a real-time macroeconomic nowcasting dashboard designed to visualize GDP growth predictions using an ensemble of econometric and machine learning models. 

Building a financial terminal requires balancing immense data density with intuitive usability. This document covers the evolution of our technology stack, our core design principles, and a thorough code-level explanation of the React-based frontend components that power the user interface. By examining the thought process behind each design choice, this report serves as both a technical manual and a UX case study.

\section{Evolution of the Technology Stack}
The development of the EconBet dashboard underwent several iterations. Our primary challenge was finding a framework that could handle complex, interconnected state changes while allowing for pixel-perfect UI customization.

\subsection{Iteration 1: Plotly Dash}
Our initial prototype was built using \textbf{Plotly Dash}. Because our backend models (Autoregressive, ADL, Random Forest) were written in Python, Dash seemed like the natural choice for seamless integration. 

\textbf{The UX Impact:} Dash allowed us to get charts on a screen quickly. However, the user experience suffered significantly. Dash applications rely heavily on server-side callbacks. Every time a user toggled a model or moved a slider, the app had to make a round-trip to the Python server. This introduced noticeable latency, making the UI feel sluggish and unresponsive.

\textbf{The Technical Bottleneck:} As the dashboard grew, we encountered "callback hell." Managing state across multiple interconnected components (e.g., changing a scenario window updating the KPI ribbon, the chart, and the ragged edge table simultaneously) resulted in convoluted, hard-to-maintain Python code. Furthermore, customizing the CSS to achieve our desired "modern financial terminal" aesthetic was clunky, as Dash abstracts away much of the underlying HTML/CSS.

\subsection{Iteration 2: Streamlit}
Seeking a simpler state management model and faster iteration cycles, we migrated the prototype to \textbf{Streamlit}.

\textbf{The UX Impact:} Streamlit provided a cleaner, more modern default look than Dash. However, the UX was rigidly constrained. Streamlit's top-down execution model meant that any interaction caused the entire script to re-run. While caching helped, the UI still flickered, and we could not create the highly interactive, localized component updates that users expect from modern web applications.

\textbf{The Technical Bottleneck:} Streamlit lacks fine-grained control over the DOM layout. Implementing a fixed sidebar with custom toggle switches, a highly customized Recharts-style tooltip, and a bespoke onboarding tour was nearly impossible without writing extensive, hacky HTML/JS injections. The UI felt too generic and did not meet our standard for a production-grade SaaS product.

\subsection{Final Decision: React + Vite + Tailwind CSS}
To achieve ultimate flexibility, rendering speed, and aesthetic control, we decoupled the frontend from the Python backend entirely and rebuilt it using \textbf{React} (bootstrapped with Vite) and \textbf{Tailwind CSS}.

\textbf{The UX Impact:} The transition to React fundamentally transformed the user experience. Interactions became instantaneous because state changes are handled entirely in the client's browser. The UI feels snappy, fluid, and professional.

\textbf{The Technical Reasoning:} 
\begin{itemize}
    \item \textbf{Component-Based Architecture:} React allowed us to build highly reusable, encapsulated UI elements.
    \item \textbf{Client-Side State:} The \texttt{useState} and \texttt{useMemo} hooks provided a clean, predictable way to manage complex application state without server round-trips.
    \item \textbf{Tailwind CSS:} Enabled us to rapidly implement a flawless Dark/Light mode toggle and achieve the exact "Bloomberg Terminal meets modern fintech" aesthetic.
    \item \textbf{Recharts:} Gave us granular, SVG-level control over our complex GDP trajectory visualizations.
\end{itemize}

\section{UI/UX Philosophy and Design System}

\subsection{The "Modern Terminal" Aesthetic}
Financial professionals are accustomed to high data density (e.g., the Bloomberg Terminal). However, traditional terminals often suffer from poor typography, harsh colors, and steep learning curves. Our design philosophy was to merge the data density of a traditional terminal with the clean, accessible aesthetics of modern consumer SaaS.

We achieved this by utilizing a strict grid system, consistent padding, and a muted color palette. We rely heavily on Tailwind CSS to enforce these design tokens across the application.

\subsection{Color Theory and Cognitive Load}
Color in EconBet is used strictly for conveying information, not for decoration. 
\begin{itemize}
    \item \textbf{Model Colors:} Each model is assigned a distinct, color-blind-friendly hue (e.g., Blue for AR, Purple for ADL, Orange for Random Forest). These colors remain consistent across the sidebar toggles, the main chart lines, and the tooltips, creating a strong mental association for the user.
    \item \textbf{Confidence Interval (Grey):} Initially, our confidence interval was colored green. User testing revealed this was distracting and drew the eye away from the actual prediction lines. We changed the confidence interval to a neutral grey (\texttt{\#808080}). This design choice pushes the uncertainty bounds into the visual background, reducing cognitive load while still providing necessary statistical context.
    \item \textbf{Directional Colors:} Green is strictly reserved for positive economic momentum, and Red for negative momentum.
\end{itemize}

\subsection{Dark Mode as a First-Class Citizen}
Financial dashboards are often viewed for hours at a time. A poorly implemented dark mode causes eye strain. We built the application to support both Light and Dark modes natively using Tailwind's \texttt{dark:} variant classes, ensuring contrast ratios meet WCAG accessibility standards in both environments.

\section{Core Architecture \& State Management}

\subsection{The Application Shell (\texttt{App.tsx})}
The core state of the application is managed at the top level in \texttt{App.tsx}. This allows data to flow downwards to child components predictably.

\begin{lstlisting}[language=JavaScript, caption=Core State Initialization]
const [data, setData] = useState<ForecastData[]>([]);
const [models, setModels] = useState({
  ar: true,
  rf: false,
  adl: false,
  ensemble: true,
  showCI: true,
});
const [useCustomWeights, setUseCustomWeights] = useState(false);
const [weights, setWeights] = useState({ ar: 33, adl: 33, rf: 34 });
const [timeRange, setTimeRange] = useState<[number, number]>([0, 0]);
const [isTourOpen, setIsTourOpen] = useState(false);
\end{lstlisting}

\textbf{Design Choice:} By lifting state to \texttt{App.tsx}, we ensure that when a user changes a setting in the \texttt{Sidebar} (like toggling a model), the \texttt{HeroChart} and \texttt{ForwardOutlook} components re-render simultaneously and synchronously.

\subsection{Performance Optimization with \texttt{useMemo}}
Parsing and manipulating large time-series datasets can cause UI stuttering if done on every render. We utilize React's \texttt{useMemo} hook heavily to cache expensive calculations.

\begin{lstlisting}[language=JavaScript, caption=Dynamic Ensemble Calculation]
const processedData = useMemo(() => {
  if (data.length === 0) return [];
  return data.map(d => {
    let ensemble = d.forecast_Ensemble;
    // Only recalculate if the user has activated custom weights
    if (useCustomWeights) {
      const total = weights.ar + weights.adl + weights.rf;
      ensemble = (d.forecast_AR * (weights.ar / total)) + 
                 (d.forecast_ADL * (weights.adl / total)) + 
                 (d.forecast_RF * (weights.rf / total));
    }
    return { ...d, forecast_Ensemble: ensemble };
  });
}, [data, useCustomWeights, weights]);
\end{lstlisting}

\textbf{UX Impact:} When a user drags the custom weight sliders, the chart updates at 60 frames per second. If we did not memoize this data, the entire component tree would re-evaluate the array mapping on every pixel the slider moved, resulting in a frustrating, laggy experience.

\section{Component Deep Dives}

\subsection{The Control Panel (\texttt{Sidebar.tsx})}
The Sidebar is the user's primary control surface. It is fixed to the left side of the screen, ensuring that controls are always accessible regardless of how far the user scrolls down the main data view.

\subsubsection{Custom Weight Sliders}
When the "Ensemble Mix" is activated, users can manually adjust the weighting of the underlying models. 

\textbf{Design Choice:} Instead of using text input fields (which require clicking, typing, and validation), we implemented range sliders using native HTML \texttt{<input type="range">}. We styled the slider thumbs using Tailwind to match the exact color of the corresponding model on the chart. This creates a direct, tactile connection between the user's hand movement and the visual output on the graph.

\subsubsection{Scenario Window Navigation}
The Scenario Window allows users to jump to specific macroeconomic events (e.g., the 2008 Financial Crisis, the COVID-19 shock). 

\textbf{UX Impact:} Navigating time-series data via a horizontal scrollbar is tedious. By providing pre-defined "Scenarios," we allow users to instantly snap the viewport to points of historical interest. When a scenario is selected, a \texttt{useEffect} hook intercepts the change and automatically adjusts the \texttt{timeRange} state, providing immediate context.

\subsection{The Main Visualization (\texttt{HeroChart.tsx})}
The centerpiece of the dashboard is the \texttt{HeroChart}, built using \textbf{Recharts}. 

\subsubsection{Visualizing Uncertainty}
Macroeconomic forecasting is inherently uncertain. Visualizing this uncertainty without overwhelming the user is a delicate balance.

\begin{lstlisting}[language=JavaScript, caption=Rendering the Confidence Interval]
{models.showCI && (
  <Area
    type="monotone"
    dataKey="ci_high"
    stroke="none"
    fill="#808080" /* Neutral Grey */
    fillOpacity={0.15}
    isAnimationActive={false}
  />
)}
\end{lstlisting}

\textbf{Design Choice:} As mentioned in the color theory section, we explicitly chose a low-opacity grey (\texttt{\#808080}) for the \texttt{Area} component representing the confidence interval. If this area were brightly colored, it would dominate the visual hierarchy. By making it grey, it acts as a subtle background element that informs the user of the variance without competing with the primary trend lines.

\subsubsection{The Custom Tooltip}
Default charting tooltips are often messy and unformatted. We engineered a custom tooltip component that intercepts the Recharts hover event. It formats the raw decimal data into clean percentages (e.g., \texttt{2.45\%}), aligns the data in a readable table, and color-codes the text to match the lines. This ensures that when a user hovers over a dense cluster of lines, they can instantly parse the exact values.

\subsection{Forward Outlook \& Trajectory (\texttt{ForwardOutlook.tsx})}
The Forward Outlook tab translates raw numerical forecasts into actionable economic narratives. 

\subsubsection{From Bar Chart to Line Chart}
Initially, we visualized the forward trajectory ($t_0$ to $t_1$ to $t_2$) using a Bar Chart. 

\textbf{Design Choice:} We refactored this to a \texttt{LineChart}. GDP growth is a continuous flow, not a set of discrete, isolated buckets. A line chart accurately conveys the \textit{momentum} and \textit{slope} of the economy. A steep downward line instantly communicates a crash much more effectively than two adjacent bars of different heights.

\subsubsection{The Economic State Logic Matrix}
To further reduce cognitive load, we implemented an algorithm that analyzes the delta (percentage point change) between time periods and categorizes the economy into specific states.

\begin{lstlisting}[language=JavaScript, caption=Economic State Logic Matrix]
function getEconomicState(prev: number, curr: number) {
  const delta = curr - prev;
  const isPositive = curr >= 0;
  const isIncreasing = delta > 0;

  if (isPositive && isIncreasing) return { 
      label: "Acceleration", desc: "Booming & gaining momentum", 
      color: "text-green-500", deltaColor: "text-green-500" 
  };
  if (isPositive && !isIncreasing) return { 
      label: "Deceleration", desc: "Growing, but cooling off", 
      color: "text-yellow-500", deltaColor: "text-red-500" 
  };
  if (!isPositive && !isIncreasing) return { 
      label: "Deepening Recession", desc: "Shrinking, decline getting worse", 
      color: "text-red-500", deltaColor: "text-red-500" 
  };
  if (!isPositive && isIncreasing) return { 
      label: "Bottoming Out", desc: "Shrinking, but crash slowing down", 
      color: "text-blue-500", deltaColor: "text-green-500" 
  };
  
  return { label: "Stable", desc: "No significant change", color: "text-gray-500" };
}
\end{lstlisting}

\textbf{UX Impact:} Users no longer have to do mental math to understand the forecast. The UI explicitly tells them: "The economy is at 2.1\%, but it is \textbf{Decelerating} (down -0.5pp from last quarter)." The color-coded labels (Green for Acceleration, Red for Deepening Recession) allow for split-second comprehension of complex econometric outputs.

\subsection{The Onboarding Engine (\texttt{OnboardingTour.tsx})}
Complex dashboards can be intimidating to new users. Initially, we used a static "Read Me" modal. We realized this was ineffective; users rarely read walls of text. We replaced it with an interactive, step-by-step Onboarding Tour.

\textbf{Design Choice:} Instead of installing a heavy third-party library like \texttt{react-joyride} (which bloats the bundle size and is difficult to style), we built a custom tour engine using \texttt{framer-motion}.

\begin{lstlisting}[language=JavaScript, caption=Dynamic Spotlight Calculation]
// Inside OnboardingTour.tsx
useEffect(() => {
  const updateSpotlight = () => {
    const target = document.querySelector(`[data-tour="${currentStepData.target}"]`);
    if (target) {
      const rect = target.getBoundingClientRect();
      setSpotlightStyle({
        top: rect.top - 8,
        left: rect.left - 8,
        width: rect.width + 16,
        height: rect.height + 16,
      });
    }
  };
  updateSpotlight();
  window.addEventListener('resize', updateSpotlight);
  return () => window.removeEventListener('resize', updateSpotlight);
}, [currentStep, currentStepData]);
\end{lstlisting}

\textbf{The Technical Implementation:} The component searches the DOM for elements tagged with specific \texttt{data-tour} attributes. It uses \texttt{getBoundingClientRect()} to calculate the exact position and dimensions of the target element on the screen. It then renders a full-screen SVG overlay with a transparent "cutout" that perfectly frames the target element, dimming the rest of the screen. 

\textbf{UX Impact:} This forces the user's attention exactly where we want it. By animating the spotlight's movement using \texttt{framer-motion}, we create a smooth, guided narrative through the application. We also implemented a \texttt{localStorage} check so the tour automatically launches for first-time visitors, but remembers if they click "Never show again."

\section{Future Improvements and Roadmap}

While the current React architecture is robust, there are several avenues for future enhancement:

\begin{itemize}
    \item \textbf{WebSockets for Live Data Streaming:} Currently, the application parses a static CSV file on load. To make this a true "Nowcasting" terminal, we plan to integrate WebSockets. This will allow the Python backend to push new GDP predictions to the React frontend in real-time as new macroeconomic indicators (like non-farm payrolls or CPI) are released, updating the charts instantly without a page refresh.
    \item \textbf{Mobile-First Responsiveness:} The current data density is optimized for desktop and tablet displays. Viewing complex multi-line charts on a mobile device is challenging. Future iterations will implement a bottom-sheet navigation pattern and simplified sparklines specifically for mobile viewports.
    \item \textbf{User Authentication \& Workspaces:} Integrating Firebase Authentication will allow users to create accounts. We can then save their custom ensemble weights, preferred scenario windows, and color themes to a database, allowing them to pick up exactly where they left off across different devices.
\end{itemize}

\section{Conclusion}

The transition from Python-based rendering frameworks (Dash, Streamlit) to a custom React architecture was a pivotal, albeit labor-intensive, decision for the EconBet project. 

By taking full control of the DOM, state management, and CSS, we transformed a clunky prototype into a highly interactive, performant financial terminal. Every design choice—from the grey confidence intervals to the custom onboarding spotlight and the economic state logic matrix—was driven by a singular goal: reducing cognitive load. We successfully abstracted the complexity of ensemble machine learning models behind an intuitive, responsive, and aesthetically polished user interface. The component-driven design ensures that the codebase remains maintainable, scalable, and ready for future real-time data integrations.

\end{document}
